\section*{Problema 1}

Encuentra una ecuación paramétrica para la recta que pasa por los siguientes puntos:

1. $(1, 1, -1)$ y $(-2, 1, 3)$

\textbf{Solución:}

Para encontrar la ecuación paramétrica de la recta que pasa por dos puntos dados, utilizamos la fórmula:

\[
\vec{r}(t) = \vec{r}_0 + t \vec{v}
\]

donde $\vec{r}_0$ es un punto en la recta y $\vec{v}$ es un vector director de la recta.

Dado que los puntos son $(1, 1, -1)$ y $(-2, 1, 3)$, podemos tomar $\vec{r}_0 = (1, 1, -1)$ y calcular el vector director $\vec{v}$ como:

\[
\vec{v} = (-2 - 1, 1 - 1, 3 - (-1)) = (-3, 0, 4)
\]

Por lo tanto, la ecuación paramétrica de la recta es:

\[
\vec{r}(t) = (1, 1, -1) + t(-3, 0, 4)
\]

O, de forma explícita:

\[
\begin{aligned}
    x &= 1 - 3t \\
    y &= 1 \\
    z &= -1 + 4t
\end{aligned}
\]

2. $(-1, 5, 2)$ y $(3, -4, 1)$

\textbf{Solución:}

De manera similar, tomamos $\vec{r}_0 = (-1, 5, 2)$ y calculamos el vector director $\vec{v}$ como:

\[
\vec{v} = (3 - (-1), -4 - 5, 1 - 2) = (4, -9, -1)
\]

Por lo tanto, la ecuación paramétrica de la recta es:

\[
\vec{r}(t) = (-1, 5, 2) + t(4, -9, -1)
\]

O, de forma explícita:

\[
\begin{aligned}
    x &= -1 + 4t \\
    y &= 5 - 9t \\
    z &= 2 - t
\end{aligned}
\]
